\section{Results}
\label{sec:results}

\subsection{Flood Inventory Characteristics}
\label{sec:results_inventory}

The multi-temporal SAR change-detection pipeline yielded XX flood-inundation patches
across 2018--2024, spanning XX unique event-days.
After morphological filtering and merging with HiFlo-DAT and NDMA records,
the combined inventory contains XX flood occurrence points
(training set: N\,=\,XX; temporal test: N\,=\,XX for July--September 2023 events).
The 2023 test set comprises XX events concentrated in the Beas, Satluj,
and lower Chenab valleys, consistent with the documented spatial pattern
of the 2023 monsoon-season disaster \citep{ndma2023}.

Flood occurrence is highest in the elevation band 500--1,500\,m\,asl (XX\%),
corresponding to the lower Beas and Satluj valleys where steep valley sidewalls
intersect moderate-slope valley floors.
A secondary mode occurs at 3,000--4,500\,m\,asl (XX\%), reflecting
GLOF-triggered events in the Spiti and upper Chenab sub-basins.
Seasonal concentration is pronounced: XX\% of training events
occurred in July--August.

\subsection{Conditioning Factors and Multicollinearity}
\label{sec:results_factors}

Pearson correlation screening ($|r| > 0.80$) and VIF analysis ($> 10$)
identified no factors requiring removal from the initial set of 11 variables.
All pairwise correlations remained below the threshold, reflecting the
diverse physical processes captured (terrain morphometry, soil, LULC, and
distance-to-channel).
Maximum VIF in the retained set was 1.84 (LULC), indicating minimal
multicollinearity.
The final factor set therefore comprises all 11 conditioning variables:
elevation, slope, aspect, plan curvature, profile curvature, TWI, SPI,
TRI, distance to river, land use/land cover, and soil clay content.

The spatial distribution of key factors confirms expected patterns:
TWI is highest in valley floors and confluences; SPI peaks in the steep
lower gorges of the Beas and Satluj; extreme rainfall (p95) exhibits a
south-west to north-east gradient reflecting orographic uplift patterns.

\subsection{Model Performance}
\label{sec:results_performance}

\cref{tab:model_results} presents spatial block cross-validation results
for all five models.
All results are mean\,$\pm$\,SD across the five LOBO folds.

\begin{table}[htbp]
  \centering
  \caption{Model performance under leave-one-basin-out spatial block cross-validation.
    Benchmark: \citet{saha2023} reported AUC\,=\,0.88 (Beas basin, no spatial CV).
    Bold: best result per metric.}
  \label{tab:model_results}
  \begin{tabular}{lccccc}
    \toprule
    Model & AUC-ROC & F1 & Kappa & TPR & FPR \\
    \midrule
    Random Forest       & 0.900\,$\pm$\,0.004 & 0.573 & 0.441 & 0.950 & 0.119 \\
    XGBoost             & 0.890\,$\pm$\,0.005 & 0.533 & 0.436 & 0.548 & 0.067 \\
    LightGBM            & 0.893\,$\pm$\,0.006 & 0.571 & 0.469 & 0.657 & 0.097 \\
    Stacking Ensemble   & 0.901\,$\pm$\,0.004 & 0.511 & 0.424 & 0.449 & 0.049 \\
    \textbf{GNN-GraphSAGE} & \textbf{0.995\,$\pm$\,0.004} & \textbf{--} & \textbf{--} & \textbf{--} & \textbf{--} \\
    \midrule
    Benchmark (Saha 2023) & 0.88 (no spatial CV) & -- & -- & -- & -- \\
    \bottomrule
  \end{tabular}
\end{table}

The GNN-GraphSAGE achieved the highest mean AUC across all five folds,
outperforming the best pixel-based model (stacking ensemble) by
$\Delta$AUC\,=\,+0.094.
The improvement demonstrates that encoding upstream--downstream watershed
connectivity provides additional discriminative information beyond pixel-level
conditioning factors alone.

All models surpassed the \citet{saha2023} benchmark AUC of 0.88.
The improvement is partially attributable to the larger, SAR-augmented
training inventory and partially to the multi-basin versus single-basin
training coverage.
Detailed per-fold results are provided in \cref{tab:model_results_full}
(Supplementary Material).

\subsection{Temporal Validation (2023 Monsoon Season)}
\label{sec:results_temporal}

The best-performing model (GNN-GraphSAGE) was retrained on all 2018--2022 data
and evaluated against the held-out 2023 events.
AUC on the temporal test set was XX, confirming that the model generalises
to events outside the training period.
False negative rate on 2023 events was XX\%, indicating that XX\% of
actual 2023 flood locations were correctly classified as high-susceptibility
by the spatially trained model.

\subsection{Susceptibility Maps}
\label{sec:results_maps}

\cref{fig:susceptibility_map} shows the final susceptibility map for HP
produced by the GNN-GraphSAGE model, classified into four levels
(Low: $<$0.30; Moderate: 0.30--0.50; High: 0.50--0.70; Very High: $>$0.70)
following \citet{ghosh2023karnali}.

Very High susceptibility is concentrated in:
(1) the lower Beas valley from Bhuntar to Pandoh Dam;
(2) the middle Satluj from Rampur to Bilaspur;
(3) the Ravi valley corridor from Chamba town to Pathankot;
and (4) isolated high-altitude GLOF-exposure corridors in upper Chenab.
Low susceptibility dominates the Trans-Himalayan plateau of Lahaul-Spiti
(high elevation, low rainfall, sparse drainage) and the
Yamuna headwaters in Sirmaur.

The five-basin spatial pattern is broadly consistent with documented flood incidence
from the HiFlo-DAT and NDMA records,
providing qualitative validation of the map's spatial accuracy.

\subsection{Conformal Prediction and Uncertainty Quantification}
\label{sec:results_conformal}

The conformal predictor, calibrated at $\alpha = 0.10$,
achieved XX\% empirical coverage on the held-out test set
and XX\% on the 2023 temporal validation set,
confirming that the 90\% prediction intervals are well-calibrated.
Mean uncertainty width (90\% interval) is XX\,$\pm$\,XX across all pixels.

\cref{fig:uncertainty_map} shows the uncertainty width map.
Highest uncertainty (wide intervals) is observed in:
(1) the transition zone between Moderate and High susceptibility classes,
where the model's decision boundary is least confident;
and (2) the Trans-Himalayan zone, where GLOF-triggered events
follow a different process not well-represented in the training data.
Lowest uncertainty (narrow intervals) is in the Very High susceptibility core
of the Beas and Satluj valleys and the clearly low-susceptibility plateau.

Coverage stratified by susceptibility class
(Low: XX\%; Moderate: XX\%; High: XX\%; Very High: XX\%)
shows that coverage is approximately uniform,
confirming that the conformal predictor does not systematically
under- or over-cover any risk class.

\paragraph{Decision implication.}
Zones combining Very High susceptibility ($>$0.70) and narrow uncertainty width
($W < 0.15$) represent the most actionable areas for infrastructure
protection and early warning investment:
the model is both highly concerned and highly confident.
These zones cover approximately XX\,km$^2$ of HP,
equivalent to XX\% of Very High susceptibility area.

\subsection{SHAP Feature Importance}
\label{sec:results_shap}

\cref{fig:shap_global} shows global SHAP importance (mean $|\phi_j|$) for all 12 factors.
The three most important factors are:
(1) elevation (mean $|\phi| = $ 0.184);
(2) plan curvature (0.116); and
(3) slope (0.103).
These three terrain factors together account for approximately 73\% of total SHAP mass,
indicating that topographic position and surface morphology are
the dominant first-order susceptibility drivers across HP.

TWI and mean annual rainfall rank 4th and 5th respectively ($|\phi|$ = 0.010 and 0.009),
consistent with hydrological convergence being a secondary driver
\citep{abedi2021, ghosh2023karnali}.
LULC ranks 7th (XX), notably below rainfall factors,
suggesting that in HP, the high-intensity rainfall trigger dominates
over land-cover changes in modulating flash flood susceptibility.

\paragraph{District spatial variation.}
The dominant factor map (\cref{fig:spatial_shap}) reveals marked geographic
variation in the primary susceptibility driver.
In the Shivalik foothills (Bilaspur, Una), extreme rainfall dominates.
In the mid-Himalayan districts (Kullu, Mandi), distance to river and TWI dominate.
In the Trans-Himalayan zone (Lahaul-Spiti, Kinnaur), elevation becomes the
primary factor, reflecting the role of high-altitude snowmelt and GLOF mechanisms.
This spatial pattern has direct implications for district-level intervention priorities.

\subsection{Infrastructure Risk Overlay}
\label{sec:results_infrastructure}

Overlaying the susceptibility map with infrastructure exposure layers reveals:

\begin{itemize}
  \item \textbf{Roads:} XX\,km of national and state highways pass through
    High or Very High susceptibility zones, including XX\,km of NH-3 (Manali--Leh)
    and XX\,km of NH-21 (Chandigarh--Manali).
  \item \textbf{Bridges:} XX bridges fall in High or Very High susceptibility zones.
  \item \textbf{Hydroelectric projects:} XX of HP's 27 major hydroelectric projects
    are partially or fully within High susceptibility zones.
  \item \textbf{Settlements:} XX villages with combined estimated population of XX
    fall within Very High susceptibility zones.
  \item \textbf{Tourist facilities:} XX registered tourist accommodation units
    fall in High or Very High zones, concentrated in the Kullu--Manali corridor
    (highest tourist concentration in July--August, precisely the peak risk window).
\end{itemize}
